\documentclass[11pt]{article}
\usepackage[a4paper,margin=1in]{geometry}
\usepackage{fontspec}
\usepackage[boldfont=false]{xeCJK}
\setCJKmainfont{FandolSong-Regular.otf}
\usepackage{amsmath}
\usepackage{amssymb}
\usepackage{array}
\usepackage{booktabs}
\usepackage{graphicx}
\usepackage{adjustbox}
\usepackage{enumitem}
\setlist[itemize]{topsep=2pt,partopsep=0pt,parsep=0pt,itemsep=2pt,leftmargin=1.4em}
\setlist[enumerate]{topsep=2pt,partopsep=0pt,parsep=0pt,itemsep=2pt,leftmargin=1.6em}
\usepackage{parskip}
\usepackage{fvextra}
\fvset{breaklines=true,breakanywhere=true,fontsize=\small}
\setlength{\parindent}{0pt}

\begin{document}

\begin{center}
  {\LARGE\bfseries Coordinate geometry}\\[0.35em]
  {\large Igcse Mathematics}
\end{center}
\vspace{0.6em}

This handout covers Topic 3, Coordinate geometry. Parts marked \textbf{(Extended)} are only tested on the Extended papers; everything else is for both levels.

\section*{Coordinates}

A point on a graph is described by its \textbf{coordinates} 坐标, sometimes called Cartesian coordinates, written $(x, y)$. The first number is the across value and the second is the up value.

\begin{itemize}
\item The two number lines are the \textbf{axes} 坐标轴: the \textbf{horizontal} 水平 $x$-axis and the \textbf{vertical} 竖直 $y$-axis.

\item They cross at the \textbf{origin} 原点, the point $(0, 0)$.

\item The axes split the grid into four \textbf{quadrants} 象限.

\end{itemize}

So the point $(3, -2)$ is found by going $3$ to the right and $2$ down.

\section*{The equation of a straight line}

Most straight lines can be written as

\[
y = mx + c,
\]

where $m$ is the \textbf{gradient} 斜率 (steepness) and $c$ is the \textbf{intercept} 截距 — the $y$-value where the line crosses the $y$-axis.

\begin{itemize}
\item A line like $x = k$ (for example $x = 3$) is \textbf{vertical}.

\item A line like $y = k$ (for example $y = 3$) is \textbf{horizontal}.

\end{itemize}

A line may also be given as $ax + by = c$. Rearrange it into $y = mx + c$ to read off the gradient and intercept.

\textbf{Worked example.} Find the gradient and $y$-intercept of $5x + 4y = 8$.

\[
4y = -5x + 8 \;\Rightarrow\; y = -\tfrac{5}{4}x + 2.
\]

So the gradient is $-\tfrac{5}{4}$ and the $y$-intercept is $2$.

\section*{Gradient}

The \textbf{gradient} measures how steep a line is:

\[
m = \frac{\text{change in } y}{\text{change in } x} = \frac{\text{rise}}{\text{run}}.
\]

A positive gradient goes up to the right; a negative gradient goes down to the right.

\textbf{Worked example.} Find the gradient of the line through $(1, 2)$ and $(4, 11)$.

\[
m = \frac{11 - 2}{4 - 1} = \frac{9}{3} = 3.
\]

\section*{Drawing a straight-line graph}

To draw $y = mx + c$, the quickest way is:

\begin{enumerate}
\item Mark the intercept $c$ on the $y$-axis.

\item Use the gradient to step to more points (for $m = 3$, go $1$ right and $3$ up).

\item Join the points with a straight line.

\end{enumerate}

You can also make a small \textbf{table of values} 数值表: choose two or three $x$-values, work out $y$, and plot the points.

\section*{Finding the equation of a line}

If you know the gradient $m$ and one point on the line, put the point into $y = mx + c$ to find $c$.

\textbf{Worked example.} A line has gradient $3$ and passes through $(1, 2)$. Find its equation.

\[
y = 3x + c, \qquad 2 = 3(1) + c, \qquad c = -1.
\]

So the equation is $y = 3x - 1$. (If you are given two points, first find the gradient, then do this.)

\section*{Length and midpoint (Extended)}

A \textbf{line segment} 线段 is the straight piece between two points.

To find the \textbf{length} 长度 of the segment between $(x_1, y_1)$ and $(x_2, y_2)$, use \textbf{Pythagoras' theorem} 勾股定理 on the horizontal and vertical gaps:

\[
\text{length} = \sqrt{(x_2 - x_1)^2 + (y_2 - y_1)^2}.
\]

To find the \textbf{midpoint} 中点 (the point exactly in the middle), average the coordinates:

\[
\text{midpoint} = \left( \frac{x_1 + x_2}{2}, \; \frac{y_1 + y_2}{2} \right).
\]

\textbf{Worked example.} Find the length and midpoint of the segment from $(1, 2)$ to $(4, 6)$.

\[
\text{length} = \sqrt{(4 - 1)^2 + (6 - 2)^2} = \sqrt{9 + 16} = \sqrt{25} = 5.
\]

\[
\text{midpoint} = \left( \frac{1 + 4}{2}, \; \frac{2 + 6}{2} \right) = (2.5, \, 4).
\]

\section*{Parallel lines}

\textbf{Parallel} 平行 lines never meet, so they have the \textbf{same gradient}.

\textbf{Worked example.} Find the equation of the line parallel to $y = 4x - 1$ that passes through $(1, -3)$.

The gradient is also $4$. Put the point in:

\[
-3 = 4(1) + c \;\Rightarrow\; c = -7,
\]

so the line is $y = 4x - 7$.

\section*{Perpendicular lines (Extended)}

Two lines are \textbf{perpendicular} 垂直 if they meet at a \textbf{right angle} 直角. Their gradients multiply to $-1$:

\[
m_1 \times m_2 = -1, \qquad \text{so} \qquad m_2 = -\frac{1}{m_1}.
\]

In words: flip the fraction and change the sign.

\textbf{Worked example.} Find the gradient of a line perpendicular to $2y = 3x + 1$.

Rearrange: $y = \tfrac{3}{2}x + \tfrac{1}{2}$, so the gradient is $\tfrac{3}{2}$. The perpendicular gradient is $-\tfrac{2}{3}$.

\subsection*{Perpendicular bisector}

The \textbf{perpendicular bisector} 垂直平分线 of a segment cuts it in half at a right angle. To find its equation: get the midpoint, then use the perpendicular gradient through that midpoint.

\textbf{Worked example.} Find the perpendicular bisector of the segment joining $(-3, 8)$ and $(9, -2)$.

\begin{itemize}
\item Midpoint: $\left( \frac{-3 + 9}{2}, \frac{8 + (-2)}{2} \right) = (3, 3)$.

\item Gradient of the segment: $\frac{-2 - 8}{9 - (-3)} = \frac{-10}{12} = -\tfrac{5}{6}$.

\item Perpendicular gradient: $\frac{6}{5}$.

\end{itemize}

Through $(3, 3)$: $\; 3 = \tfrac{6}{5}(3) + c \Rightarrow c = 3 - \tfrac{18}{5} = -\tfrac{3}{5}$. So the bisector is

\[
y = \tfrac{6}{5}x - \tfrac{3}{5}.
\]


\end{document}
